% Vietnamese translation for OI Public Library
% Source-PDF: IOI中国国家候选队论文/2000/郭一论文.pdf
\documentclass[11pt]{article}
\usepackage{oipl-vn}

\newcommand{\Cat}{\operatorname{Cat}}

\title{Mô Hình Toán Học Và Ứng Dụng Trong Thi Tin Học}
\author{Guo Yi, Trường Trung học trực thuộc Đại học Phục Đán}
\date{2000}

\begin{document}
\maketitle

\origtitle{Mathematical models and their applications in informatics contests}
\sourcepath{IOI Candidate Team Papers/2000/Guo Yi.pdf}

\begin{abstract}
Mô hình toán học là một công cụ mạnh để giải các bài toán thực tế. Trong thi
tin học, điều thường gặp không phải là mô hình liên tục đòi hỏi nhiều kiến
thức giải tích, mà là các mô hình rời rạc. Bài viết giới thiệu khái niệm chung
của mô hình toán học rời rạc, tiêu chuẩn lựa chọn mô hình, và vai trò của việc
chuyển hóa bài toán thực tế thành mô hình quen thuộc.
\end{abstract}

\noindent\textbf{Từ khóa:} mô hình toán học, độ tin cậy, khả năng giải được.

\section{Mở đầu}

Mô hình toán học là vũ khí quan trọng để con người giải quyết các vấn đề trong
thực tế. Ta trừu tượng hóa và tinh luyện một vấn đề thực thành một mô hình
toán học, rồi dùng phương pháp toán học để xử lý nó. Mô hình toán học có thể
chia thành hai loại lớn: liên tục và rời rạc. Mô hình liên tục thường cần
nhiều kiến thức toán cao cấp, còn trong các kỳ thi tin học, mô hình thường gặp
là mô hình toán học rời rạc. Bài viết này tập trung vào khái niệm và phương
pháp xây dựng các mô hình như vậy.

Trong tin học thi đấu, đề bài thường mô tả một vấn đề thực tế. Sau khi đọc
hiểu đề, thí sinh cần rút ra các yếu tố then chốt và tạo thành một mô hình
trừu tượng. Mô hình tốt làm cho việc phân tích và giải bài toán trở nên rõ
ràng hơn.

Thông thường, khi giải một bài toán xuất phát từ thực tế, ta đi qua các bước:
\begin{enumerate}
  \item Đọc kỹ đề, hiểu bối cảnh, xác định dữ liệu đã biết, kết quả cần xuất
  ra, và quy mô dữ liệu.
  \item Xây dựng mô hình sao cho mô hình biểu diễn vấn đề thực một cách ngắn
  gọn và hiệu quả.
  \item Giải mô hình, từ đó thu được thuật toán. Bước này phụ thuộc rất lớn
  vào việc mô hình được xây dựng có dễ giải hay không.
  \item Cài đặt chương trình.
\end{enumerate}
Trong đó, bước thứ hai và thứ ba là then chốt. Chất lượng của mô hình quyết
định trực tiếp khả năng giải đúng và giải hiệu quả bài toán.

\section{Một vấn đề có thể có nhiều mô hình}

Hiện tượng một bài có nhiều cách giải rất phổ biến. Vì vậy ta cần bàn đến tiêu
chuẩn chọn mô hình. Một mô hình tốt không chỉ phải phản ánh chính xác mô hình
thực tế, tức có \term{độ tin cậy}; nó còn phải có thể được giải một cách hiệu
quả, tức có \term{khả năng giải được}. Ở đây ``hiệu quả'' nghĩa là thuật toán
giải mô hình có chi phí thời gian và bộ nhớ chấp nhận được.

Khi giải các bài yêu cầu tối ưu hoặc đếm chính xác, ta thường bảo đảm độ tin
cậy trước, sau đó cố gắng nâng cao khả năng giải được. Nếu nhiều mô hình cùng
đáng tin cậy, mô hình nào dễ giải hơn thường là mô hình tốt hơn.

\subsection{Ví dụ: tuyến du lịch tốt nhất}

Xét bài \textit{A better tour route} của IOI 1993. Một hãng hàng không tặng vé
miễn phí để du lịch ở Canada. Người đi xuất phát từ thành phố cực tây, bay một
chiều từ tây sang đông qua một số thành phố để tới thành phố cực đông, rồi bay
một chiều từ đông sang tây trở về điểm xuất phát. Ngoài thành phố xuất phát,
mỗi thành phố chỉ được thăm nhiều nhất một lần; thành phố xuất phát được thăm
hai lần. Chỉ được dùng các đường bay trực tiếp đã cho. Cần tìm một tuyến thỏa
các ràng buộc và thăm được nhiều thành phố nhất.

Ta có thể trừu tượng hóa mỗi thành phố thành một đỉnh. Đánh số các thành phố
từ tây sang đông là \(1,2,\ldots,n\). Nếu hai thành phố có đường bay trực tiếp,
ta nối hai đỉnh tương ứng bằng một cạnh vô hướng.

\paragraph{Tìm kiếm mù.}
Ý tưởng trực tiếp nhất là giả sử có hai máy bay cùng bay từ thành phố cực tây
đến thành phố cực đông. Trong quá trình tìm kiếm, ta bảo đảm hai tuyến không
có thành phố chung ngoài hai đầu mút. Với mọi cặp tuyến tìm được, ghi nhận cặp
đi qua nhiều thành phố nhất; một tuyến dùng cho lượt đi, tuyến kia dùng cho
lượt về.

Mô hình này đúng, nhưng độ phức tạp tìm kiếm là cấp số mũ. Nguyên nhân chính
là mức trừu tượng còn thấp: các ràng buộc của đề chưa được đưa đủ vào mô hình,
nên quá trình tìm kiếm còn quá mù quáng.

\paragraph{Mô hình luồng.}
Ta có thể đưa nhiều ràng buộc hơn vào mô hình bằng cách xây dựng một bài toán
luồng cực đại chi phí cực đại:
\begin{enumerate}
  \item Để bảo đảm mỗi thành phố được thăm nhiều nhất một lần, tách mỗi thành
  phố \(i\) thành hai đỉnh \(i\) và \(i'\), rồi nối cung \(i\to i'\) có chi
  phí đơn vị bằng \(0\).
  \item Với mỗi cạnh vô hướng giữa \(i\) và \(j\), giả sử \(i<j\), thêm cung
  \(i'\to j\) có dung lượng \(1\). Cung này biểu diễn việc bay từ tây sang
  đông và chỉ dùng một lần. Chi phí của cung đặt là \(j-i+1\), biểu diễn số
  thành phố đi qua trên đoạn đó, tính cả hai đầu.
  \item Các cung \(1\to 1'\) và \(n\to n'\) có dung lượng \(2\).
\end{enumerate}

Khi đó ta tìm luồng cực đại chi phí cực đại từ \(1\) tới \(n'\). Nếu lưu lượng
cực đại không bằng \(2\), không có cách đi từ tây sang đông rồi trở về theo
đúng yêu cầu. Ngược lại, giả sử chi phí cực đại là \(C\). Vì có hai tuyến, mỗi
thành phố không được thăm đóng góp \(2\) trong cách tính chi phí, còn mỗi
thành phố được thăm đóng góp \(3\); hai thành phố đầu cuối chỉ đóng góp \(2\).
Số thành phố được thăm là \(C+2-2n\). Mô hình này đáng tin cậy và có thể giải
bằng thuật toán Ford--Fulkerson cho biến thể chi phí cực đại. So với tìm kiếm
mù, độ phức tạp giảm từ cấp số mũ xuống đa thức, xấp xỉ \(O(n^3)\).

\paragraph{Mô hình quy hoạch động.}
Ta còn có thể xây dựng mô hình gọn hơn. Gọi \(f[i,j]\) là tổng số đường bay
lớn nhất mà hai tuyến rời nhau, cùng xuất phát từ đỉnh \(1\), có thể đạt được
khi lần lượt dừng ở \(i\) và \(j\). Nếu \(i>j\), hoặc \(i=j=n\), ta xét đỉnh
trước \(i\), ký hiệu là \(k\). Khi đó trạng thái tới \(k\) và \(j\) cũng phải
tối ưu; nếu không sẽ mâu thuẫn với tính tối ưu của \(f[i,j]\). Trường hợp
\(i<j\) tương tự. Công thức có dạng:
\[
f[1,1]=0,
\]
\[
f[i,j]=
\begin{cases}
\displaystyle \max_{\substack{k<i\\(k,i)\in E}} f[k,j]+1,
  & i>j \text{ hoặc } i=j=n,\\[1.2em]
\displaystyle \max_{\substack{k<j\\(k,j)\in E}} f[i,k]+1,
  & i<j.
\end{cases}
\]
Trạng thái \(f[i,i]\) không có nghĩa với \(1<i<n\). Đáp án là
\(f[n,n]-2\), và độ phức tạp giảm xuống \(O(n^2)\).

Ba mô hình trên đều có thể cho lời giải đúng, nhưng khả năng giải được khác
nhau. Nói chung, mô hình càng ngắn gọn và mức trừu tượng càng cao thì khi cài
đặt càng ít thao tác thừa, và chương trình càng hiệu quả.

Một điểm cần chú ý là gần đây trong một số cuộc thi xuất hiện bài chấm điểm
theo mức độ xấp xỉ của lời giải. Với loại bài này, ta thường cân nhắc nhiều
hơn đến khả năng giải được của mô hình. Dĩ nhiên độ tin cậy chỉ được hy sinh
trong giới hạn cho phép: phương án vẫn phải khả thi và sai số phải nằm trong
phạm vi được chấp nhận.

\section{Một mô hình có thể áp dụng cho nhiều vấn đề}

Đây là một lý do quan trọng khiến ta nghiên cứu mô hình toán học. Giải được
một mô hình nghĩa là ta giải được cả một lớp bài toán. Chẳng hạn, bài toán
đường đi ngắn nhất xuất hiện khắp nơi trong thực tế; mô hình luồng cũng có giá
trị ứng dụng rất cao. Tuy nhiên, giải mô hình mới chỉ là một nửa công việc.
Nửa còn lại là chuyển bài toán thực tế thành một mô hình đã biết. Trong thi
đấu, nửa sau nhiều khi còn quan trọng hơn.

\section{Làm thế nào để xây dựng mô hình}

Muốn xây dựng mô hình, trước hết cần quen thuộc với một số mô hình kinh điển
và thuật toán của chúng. Đây là nền móng của việc mô hình hóa; thiếu nền móng
này thì không thể nói đến việc xây dựng mô hình. Rất nhiều bài thi cuối cùng
có thể chuyển về các mô hình quen thuộc.

Tiếp theo, thí sinh cần năng lực liên hệ và so sánh các bài toán thực tế. Như
đã nói, nhiều vấn đề thực tế có cùng một mô hình toán học hoặc mô hình rất
giống nhau. Liên hệ và so sánh là cách hữu hiệu để phát hiện điều đó.

\subsection{Các bài toán Catalan}

\paragraph{Chia đa giác lồi.}
Một đa giác lồi \(n\) cạnh có thể được chia thành \(n-2\) tam giác bằng
\(n-3\) đường chéo không cắt nhau. Hỏi số phương án là \(h_n\). Khi \(n=5\),
ta có \(h_5=5\).

\paragraph{Dãy ngoặc đúng.}
Với \(n\) dấu ngoặc trái và \(n\) dấu ngoặc phải, hỏi có bao nhiêu dãy ngoặc
đúng khác nhau, ký hiệu \(Q_n\). Quy tắc định nghĩa dãy đúng là:
\begin{enumerate}
  \item \texttt{()} là đúng.
  \item Nếu \(A\) đúng thì \texttt{(}\(A\)\texttt{)} cũng đúng.
  \item Nếu \(A\) và \(B\) đều đúng thì \(AB\) cũng đúng.
\end{enumerate}

Hai bài toán trên đều dẫn đến số Catalan:
\[
  \Cat_n=\frac{1}{n+1}\binom{2n}{n}.
\]
Cụ thể, \(h_n=\Cat_{n-2}\) và \(Q_n=\Cat_n\).

\paragraph{Luật kết hợp.}
Cho \(n\) số \(a_1,a_2,\ldots,a_n\). Giữ nguyên thứ tự và chỉ dùng ngoặc để
chỉ ra cách cộng theo luật kết hợp, hỏi có bao nhiêu phương án \(P_n\). Vì bài
toán chỉ đếm cách ghép cặp, không phụ thuộc vào giá trị cụ thể của các
\(a_i\), ta có thể hình dung các \(a_i\) là các vectơ liên tiếp. Mỗi cách đặt
ngoặc tương ứng một cách chia đa giác, do đó
\[
  P_n=h_{n+1}=\Cat_{n-1}.
\]

\paragraph{Đếm cây nhị phân.}
Hỏi có bao nhiêu cây nhị phân gồm \(n\) nút. Bài này có thể liên hệ với dãy
ngoặc đúng. Xét quá trình duyệt trung thứ tự bắt đầu từ gốc: lần đầu đến một
nút thì đưa nút đó vào ngăn xếp, ghi một dấu \texttt{(}; sau khi duyệt xong
cây con trái thì thăm nút, lấy nút ra khỏi ngăn xếp, ghi một dấu \texttt{)};
rồi tiếp tục duyệt cây con phải. Như vậy mỗi cây nhị phân tương ứng với một
dãy ngoặc đúng, và ngược lại. Mỗi nút vào ra ngăn xếp một lần, nên có \(n\)
cặp ngoặc. Số cây là \(Q_n=\Cat_n\).

Các bài tương tự còn có bài xếp bát, bài xếp hàng đổi tiền, v.v. Bề ngoài
chúng khác nhau, nhưng bản chất đều là mô hình Catalan. Liên hệ và so sánh là
cách tìm ra bản chất chung đó.

\section{Sáng tạo trong mô hình hóa}

Thi tin học không chỉ là cuộc thi lập trình, mà còn là cuộc thi trí tuệ. Sức
sáng tạo là một phần quan trọng của trí tuệ. Mô hình hóa toán học không có
khuôn mẫu cố định, vì vậy nó thể hiện rõ sức sáng tạo của thí sinh và ngày
càng được các cuộc thi tin học coi trọng. Khi đề thi phát triển, mô hình toán
học trong đề ngày càng ẩn, và yêu cầu về sức sáng tạo cũng tăng lên.

\subsection{Ví dụ: xâu 01, NOI 1999}

Bài toán yêu cầu tìm một xâu độ dài \(N\) chỉ gồm \(0\) và \(1\), sao cho:
\begin{itemize}
  \item trong mọi đoạn liên tiếp độ dài \(L_0\), số ký tự \(0\) không nhỏ hơn
  \(A_0\) và không lớn hơn \(B_0\);
  \item trong mọi đoạn liên tiếp độ dài \(L_1\), số ký tự \(1\) không nhỏ hơn
  \(A_1\) và không lớn hơn \(B_1\).
\end{itemize}
Nếu không tồn tại xâu như vậy thì in \(-1\), ngược lại in một xâu bất kỳ. Giới
hạn là \(N\le 1000\).

Mô hình của bài này khá kín đáo. Nhiều thí sinh khi đó dùng tìm kiếm mù. Nhưng
xâu dài tới \(1000\), mỗi vị trí có hai khả năng, nên độ phức tạp xấu nhất
tiệm cận \(2^{1000}\), hoàn toàn không thể chấp nhận.

Trước hết xét bài \textit{Black and White} của CEOI 1994: tìm một dãy gồm
\(n\) số nguyên sao cho tổng của mọi đoạn liên tiếp \(p\) số là dương, còn
tổng của mọi đoạn liên tiếp \(q\) số là âm. Nếu không tồn tại thì in
\texttt{NO}, ngược lại in một dãy.

Điểm then chốt là biểu diễn tổng đoạn liên tiếp. Nếu viết trực tiếp
\(a_{i+1}+a_{i+2}+\cdots+a_{i+k}\), mô hình rất khó xây dựng vì có nhiều biến
tham gia. Nhưng nếu đặt
\[
  S_i=a_1+a_2+\cdots+a_i,\qquad S_0=0,
\]
thì tổng đoạn trở thành \(S_{i+k}-S_i\). Điều kiện của bài CEOI trở thành:
\[
  S_i>S_{i-p}\quad(p\le i\le n),\qquad
  S_i<S_{i-q}\quad(q\le i\le n).
\]
Ta dựng đồ thị có \(n+1\) đỉnh \(S_0,S_1,\ldots,S_n\). Nếu cần \(S_i<S_j\),
thêm cạnh có hướng từ \(S_i\) tới \(S_j\). Khi đó dãy \(S_0,\ldots,S_n\) phải
tuân theo một thứ tự topo. Ngược lại, bất kỳ thứ tự topo phù hợp nào cũng xác
định một dãy số nguyên. Bài toán được chuyển thành topo sort.

Quay lại bài NOI 1999, ta cũng đặt \(S_i\) là số lượng ký tự \(1\) trong đoạn
từ vị trí \(1\) đến \(i\). Không chỉ có quan hệ lớn bé giữa \(S_i\) và \(S_j\),
ta còn cần biết lớn hơn bao nhiêu; vì vậy dùng cạnh có trọng số. Nếu
\(S_j\ge S_i+k\), thêm cạnh \(S_i\to S_j\) trọng số \(k\). Từ các điều kiện về
số lượng \(0\) và \(1\), ta nhận được các ràng buộc hiệu. Ngoài ra vì mỗi ký
tự chỉ là \(0\) hoặc \(1\), cần thêm:
\[
  S_{i+1}\ge S_i,\qquad S_{i+1}\le S_i+1.
\]
Sau đó bài toán trở thành tìm đường dài nhất từ \(S_0\) tới các đỉnh còn lại
trong hệ ràng buộc hiệu. Bài CEOI phía trên cũng có thể xem là một trường hợp
của cùng mô hình đồ thị đường dài nhất, với trọng số cạnh bằng \(1\).

Ban đầu hai bài toán này dường như chẳng liên quan đến đồ thị: đề không nói
đến thành phố, trạm, đường, tuyến, độ dài hay thời gian truyền. Nhưng mô hình
đồ thị lại giải chúng rất đẹp. Việc nghĩ ra các đỉnh \(S_i\), dùng quan hệ lớn
bé làm cạnh, và xác định trọng số chính là phần sáng tạo của mô hình hóa.

\section{Phụ lục: hai thuật toán trong bài gốc}

PDF gốc kèm hai chương trình Pascal. Dưới đây là bản diễn giải thuật toán,
tránh giữ lại các chú thích tiếng Trung trong mã nguồn gốc.

\subsection{Black and White, CEOI 1994}

Với các đỉnh \(S_0,\ldots,S_n\), tính bậc vào của đồ thị ràng buộc:
\begin{itemize}
  \item nếu \(i+p\le n\), thêm ràng buộc \(S_i<S_{i+p}\);
  \item nếu \(i-q\ge 0\), thêm ràng buộc \(S_i<S_{i-q}\).
\end{itemize}
Sau đó thực hiện sắp xếp topo. Mỗi khi lấy một đỉnh bậc vào \(0\), gán cho nó
thứ tự hiện tại rồi xóa các cạnh đi ra. Nếu không thể gán đủ \(n+1\) đỉnh,
đồ thị có chu trình và không có lời giải. Ngược lại, dãy cần xuất là
\[
  s_i-s_{i-1}\qquad(1\le i\le n),
\]
trong đó \(s_i\) là thứ tự topo của đỉnh \(S_i\).

\subsection{Xâu 01, NOI 1999}

Dựng đồ thị trọng số trên các đỉnh \(0,1,\ldots,n\). Với mỗi vị trí \(i\):
\begin{itemize}
  \item từ ràng buộc về số lượng \(0\) trong đoạn độ dài \(L_0\), thêm các
  cạnh tương ứng với
  \(A_0\le L_0-(S_i-S_{i-L_0})\le B_0\);
  \item từ ràng buộc về số lượng \(1\) trong đoạn độ dài \(L_1\), thêm các
  cạnh tương ứng với
  \(A_1\le S_i-S_{i-L_1}\le B_1\);
  \item thêm cạnh biểu diễn \(S_i\ge S_{i-1}\) và \(S_i\le S_{i-1}+1\).
\end{itemize}
Chạy phép thư giãn đường dài nhất từ đỉnh \(0\). Nếu sau \(n\) vòng vẫn còn
cạnh cải thiện được, hoặc có \(S_i>i\), hệ ràng buộc vô nghiệm và in \(-1\).
Nếu không, xuất ký tự thứ \(i\) là \(0\) khi \(S_i=S_{i-1}\), và là \(1\) khi
\(S_i=S_{i-1}+1\).

\section{Kết luận}

Nhìn lại lịch sử phát triển của tin học, từ phần cứng, phần mềm đến mạng máy
tính, ở đâu cũng thấy vai trò của sáng tạo. Sáng tạo là phẩm chất cơ bản của
người học tin học và là trọng điểm đánh giá của thi tin học. Mô hình toán học
giúp ta nối bài toán thực tế với các cấu trúc trừu tượng đã hiểu rõ; xây dựng
được cây cầu ấy thường là bước quyết định của lời giải.

\section*{Tài liệu tham khảo}

\begin{itemize}
  \item Wu Wenhu và Wang Jiande, \textit{Thuật toán và lập trình tổ hợp toán
  học}, bộ sách hướng dẫn Olympic Tin học thanh thiếu niên quốc tế và toàn
  quốc, Nhà xuất bản Đại học Thanh Hoa.
  \item Wu Wenhu và Wang Jiande, \textit{Thuật toán và lập trình đồ thị}, cùng
  bộ sách trên, Nhà xuất bản Đại học Thanh Hoa.
  \item Wang Shuhe, \textit{Cơ sở mô hình toán học}, Nhà xuất bản Đại học Khoa
  học và Công nghệ Trung Quốc.
\end{itemize}

\end{document}
